\section{Conclusão}

Este trabalho investigou como a composição de uma mistura de óleo de coco babaçu e óleo de soja e o uso de aditivos influenciam um sabão produzido por saponificação a quente. Confirmou-se que o perfil de ácidos graxos determina as propriedades do produto: o babaçu confere dureza e espuma, e a soja, maciez. A composição de cerca de 33\% de babaçu e 66\% de soja foi a mais equilibrada; o ácido esteárico aumentou a dureza e o etanol foi o aditivo de grande impacto, gerando a melhor amostra (quinto dia). Após a cura, os produtos atingiram pH em torno de 9, espuma estável e boa capacidade de limpeza, cumprindo os objetivos do trabalho.

O principal passo, que não foi tomado, seria esverdear a rota, aprofundando seu alinhamento aos princípios da Química Verde. A modificação de maior impacto seria a substituição, parcial ou total, do óleo de soja por óleo residual de fritura ou semelhantes, evitando o descarte inadequado desse resíduo; isso exigiria uma etapa prévia de filtração e o reajuste do índice de saponificação, em razão da variabilidade na composição do óleo reutilizado. Somam-se a essa medida a redução do excesso de NaOH, que diminui tanto o consumo de reagente quanto a base residual no produto, também seria interessante trabalhar na diminuição do consumo energético, por meio de temperaturas mais baixas ou de agitação manual, e o uso de materiais recicláveis no lugar dos plásticos descartáveis usados durante os experimentos.
